\documentclass[12pt]{article}
\usepackage{amsmath}
\usepackage{amssymb}
\usepackage[margin=1in]{geometry}

\title{Can a simple model explain economics?}
\author{Ludovico Furlanetto}
\date{July 2026}

\begin{document}
\maketitle

\begin{abstract}
Across datasets, the volatility of a firm's growth rate falls slowly with size,
$\sigma(S)\sim S^{-\beta}$ with $\beta\approx0.15$--$0.20$, and its growth-rate distribution is
fat-tailed. The standard explanation treats each firm as a collection of independently fluctuating
sub-units with heavy-tailed sizes, so that both facts arise from the firm's internal
composition. Moran, Secchi, and Bouchaud find that firm data contradict this granular picture, and suggest
that the shocks hitting a firm's sub-units must become more correlated as it
grows. We ask instead
whether a dynamical model of interacting firms can generate both on its own.

We model the economy as a Generalized Lotka--Volterra system on a scale-free random graph, in
which the interactions act on a firm's size relative to the average firm to
make the dynamics scale-invariant. In the strongly disordered
regime the economy's composition keeps fluctuating chaotically, and in this persistently
fluctuating state the two facts emerge together as stationary properties of the
interactions, with no external noise. On the scale-free graph the normalized higher-moment
exponents $\zeta_q/\zeta_1\approx1,\,1.84,\,2.51,\,3.02$ reproduce the anomalous multiscaling
of the empirical $1,\,2.0,\,2.6,\,2.9$, and the size--variance relation decreases with firm
size, as in the data. Its exponent is set by the graph's degree tail, which brackets the
empirical $0.15$--$0.20$: $\beta \approx 0.45$ at tail exponent $2.5$, where the ratios sit
on the data, falling to $\beta = 0.21$ at exponent $2.2$, where it reaches the band while the
ratios run a few percent low. The other facts Moran, Secchi, and Bouchaud test against the
granular picture, a fat-tailed growth-rate distribution and heavy tails that strengthen
with firm size, hold in the model as well.
\end{abstract}

\end{document}
